% SGA 3, Expose VI_A, Loop 1 English reconstruction.
% Scope: paragraph 2.5 through the end of Proposition 2.5.2;
% Lemma 2.5.3 is excluded.
% Authority: Polo--Gille Exp6A-13oct24.pdf, component-local pages 11--13
% (combined-reader pages 314--316; printed pages 303--305).
% OCR and the Reinhold English Markdown are comparison witnesses only.

\medskip
\phantomsection\label{sec:sga3-VIA-2.5}%
\noindent\SGAInlineRefTarget{sga3:vi:paragraph:2-5}\textbf{2.5.}
此段为新增内容. 以下结果出现于第 VI\textsubscript{B} 篇,
但本可 (或本应) 出现于 VI\textsubscript{A};
为了更精确地表述下文\hyperref[thm:sga3-VIA-3.2]{定理 3.2},
提前具备这些结果是有益的.

\medskip
\phantomsection\label{lem:sga3-VIA-2.5.1}%
\noindent\textbf{引理 2.5.1.}
\textit{设 \((A,\mathfrak m)\) 是 Artin 局部环,
\(k=A/\mathfrak m\) 是其剩余域.}
\begin{enumerate}
\item[\SGAInlineRefTarget{lem:sga3-VIA-2.5.1:item:i}(i)]
\textit{若 \(X\) 是一个 \(A\)-概形,
使得 \(X\otimes_Ak\) 在 \(k\) 上局部有限型 (分别地, 有限型),
则 \(X\) 在 \(A\) 上具有相同的性质.}
\item[\SGAInlineRefTarget{lem:sga3-VIA-2.5.1:item:ii}(ii)]
\textit{设 \(u:X\to Y\) 是 \(A\)-概形态射.
若 \(u\otimes_Ak\) 是浸入 (分别地, 闭浸入), 则 \(u\) 也如此.}
\end{enumerate}

\noindent\textit{证明.}
\SGARefLink{lem:sga3-VIA-2.5.1:item:i}{(i)}
假设 \(X\otimes_Ak\) 在 \(k\) 上局部有限型.
设 \(U=\operatorname{Spec}B\) 是 \(X\) 的一个仿射开子集.
由假设, \(B\) 中存在元素 \(x_1,\ldots,x_n\),
其像生成作为 \(k\)-代数的 \(B/\mathfrak mB\);
于是 ``幂零 Nakayama 引理'' 表明,
这些 \(x_i\) 生成作为 \(A\)-代数的 \(B\).
这证明 \(X\) 在 \(A\) 上局部有限型.
此外, 若 \(X\otimes_Ak\) 拟紧,
则与它有相同底拓扑空间的 \(X\) 也拟紧;
所以 \(X\) 在 \(A\) 上有限型.
这就证明了 \SGARefLink{lem:sga3-VIA-2.5.1:item:i}{(i)}.

证明 \SGARefLink{lem:sga3-VIA-2.5.1:item:ii}{(ii)}.
假设 \(u\otimes_Ak\) 是浸入 (分别地, 闭浸入).
则 \(u\) 是从 \(X\) 到 \(Y\) 的一个局部闭子集 (分别地, 闭子集) 的同胚,
并且对每个 \(x\in X\), 环同态
\[
  \varphi_x:\mathcal O_{Y,u(x)}\longrightarrow\mathcal O_{X,x}
\]
满足 \(\varphi_x\otimes_Ak\) 为满射.
因此幂零 Nakayama 引理表明 \(\varphi_x\) 是满射,
所以 \(u\) 是浸入 (分别地, 闭浸入).

\medskip
\phantomsection\label{prop:sga3-VIA-2.5.2}%
\noindent\textbf{命题 2.5.2.}
\textit{设 \(A\) 是剩余域为 \(k\) 的 Artin 局部环,
\(u:G\to H\) 是局部有限型的 \(A\)-群之间的一个拟紧态射.}
\begin{enumerate}
\item[\SGAInlineRefTarget{prop:sga3-VIA-2.5.2:item:a}(a)]
\textit{集合 \(u(G)\) 在 \(H\) 中是闭的,
且其连通分支都不可约, 维数全都相同.}
\item[\SGAInlineRefTarget{prop:sga3-VIA-2.5.2:item:b}(b)] \textit{有}
\[
  \dim G=\dim u(G)+\dim\operatorname{Ker}(u).
\]
\item[(c)] \textit{若 \(u\) 是单态射, 则它是闭浸入.}
\end{enumerate}

\noindent\textit{证明.}
由\hyperref[lem:sga3-VIA-2.5.1]{前一引理},
只需在 \(A=k\) 时证明此命题.
此外, 所讨论的性质在 fpqc 下降下稳定,
且维数在基域扩张下不变,
所以可以假设 \(k\) 代数闭.

证明 \SGARefLink{prop:sga3-VIA-2.5.2:item:a}{(a)}.
设 \(C\) 是 \(H\) 中以 \(u(G)\) 为底拓扑空间的约化子概形.
由于 \(u(G)\) 在 \(H\) 的求逆态射下稳定, \(C\) 也如此.
此外, \(u:G\to C\) 拟紧且占优.
因此, 由 EGA~IV\textsubscript{2}, 2.3.7,
\[
  u\times_k\operatorname{id}_G
  \qquad\text{及}\qquad
  \operatorname{id}_H\times_ku,
\]
也如此, 从而它们的复合
\[
  u\times_ku:G\times_kG\longrightarrow C\times_kC.
\]
也如此. 所以 \(H\) 上的乘法将 \(C\times_kC\) 映入 \(C\),
因而 \(C\) 是 \(H\) 的一个群子概形.

以 \(C\) 代替 \(H\), 便归约到 \(u\) 占优的情形.
此时 \(G(k)\) 在 \(H\) 中稠密,
因而与 \(H\) 的每个连通分支相交;
所以它在这些连通分支所成的集合上的作用是传递的.
因此只需证明 \(u(G)\) 包含 \(H^0\).
以 \(u^{-1}(H^0)\) 代替 \(G\), 便可假设 \(H=H^0\).
此时, \hyperref[prop:sga3-VIA-2.4]{2.4} 表明 \(H\) 不可约,
且在 \(k\) 上有限型, 因而为 Noether 型.
此外, \(u\) 局部有限型 (参见 EGA~I, 6.6.6) 且拟紧,
所以有限型. 由 Chevalley 可构造性定理
(参见 EGA~IV\textsubscript{1}, 1.8.5),
\(u(G)\) 因而是
\[
  H=\overline{u(G)},
\]
的可构造 (且稠密) 子集,
所以包含 \(H\) 的一个稠密开子集 \(U\)
(参见 EGA~0\textsubscript{III}, 9.2.2).
由 \hyperref[sga3:VIA:0.5]{0.5}, 于是有
\[
  H=U\cdot U\subset u(G),
\]
从而 \(u(G)=H\).
结合 \hyperref[cor:sga3-VIA-2.4.1]{2.4.1},
便证明了 \SGARefLink{prop:sga3-VIA-2.5.2:item:a}{(a)}.

证明 \SGARefLink{prop:sga3-VIA-2.5.2:item:b}{(b)}.
首先回顾, 函子 \(\operatorname{Ker}(u)\)
(参见 I, \SGARefLink{sga3:i:definition:2-3-6-1}{2.3.6.1})
由 \(u^{-1}(e)\) 表出, 其中 \(e\) 表示 \(H\) 的单位元.
由于 \(u\) 有限型, \(\operatorname{Ker}(u)\) 在 \(k\) 上有限型.
另一方面, 用约化闭子概形 \(u(G)\) 代替 \(H\),
可以假设 \(u\) 为满射.
以 \(u^0\) 表示 \(u\) 在 \(G^0\) 上的限制.
由于 \(G\) 与 \(\operatorname{Ker}(u)\) 都等维,
且
\[
  \operatorname{Ker}(u)^0\subset\operatorname{Ker}(u^0),
\]
便归约到 \(G\), 因而也有 \(H\), 不可约的情形.

由 EGA~IV\textsubscript{3}, 9.2.6.2 和 10.6.1(ii),
满足
\[
  \dim u^{-1}(y)=\dim G-\dim H
\]
的那些 \(y\in H\) 所成的集合包含一个非空开子集 \(V\).
由于 \(u\) 为满射, \(U=u^{-1}(V)\) 此时是 \(G\) 的一个非空开子集,
从而包含 \(G\) 的一个闭点 \(x\),
因为 \(G\) 是 Jacobson 概形 (参见 EGA~IV\textsubscript{3}, 10.4.8).
右平移 \(r_x\) 是从 \(\operatorname{Ker}(u)\) 到
\(u^{-1}(u(x))\) 的一个同构, 从而
\[
  \dim\operatorname{Ker}(u)
    =\dim u^{-1}(u(x))
    =\dim G-\dim H.
\]

遵循 \textup{[DG70]}, I, \S~3.4 来证明 (c).
(另一证明见第 VI\textsubscript{B} 篇, 1.4.2.)
假设 \(u\) 是单态射.
若 \(C\) 是 \(G\) 的一个连通分支,
则存在一个闭点 \(x\in G\) 使得 \(C=r_x(G^0)\).
若 \(u_C\) (分别地, \(u^0\)) 表示 \(u\) 在 \(C\)
(分别地, \(G^0\)) 上的限制, 则
\[
  u_C=r_{u(x)}\circ u^0\circ r_x^{-1}.
\]
所以只需证明 \(u^0\) 是闭浸入.
因而可以假设 \(G=G^0\),
于是 \(G\) 不可约且在 \(k\) 上有限型.

设 \(\xi\) 是 \(G\) 的泛点.
则 \(\mathcal O_{G,\xi}\) 是 Artin 局部环;
以 \(\mathfrak m\) 表示其极大理想.
设 \(h=u(\xi)\), 以 \(\mathfrak n\) 表示
\(\mathcal O_{H,h}\) 的极大理想, 并置
\[
  A=\mathcal O_{G,\xi}/\mathfrak n\mathcal O_{G,\xi}.
\]
由于 \(u\) 是单态射, 经基变换得到的态射
\[
  u_h:\operatorname{Spec}(A)\longrightarrow
      \operatorname{Spec}(\kappa(h))
\]
也是单态射. 因此乘法态射
\[
  A\otimes_{\kappa(h)}A\longrightarrow A
\]
是同构 (参见 EGA~I, 5.3.8), 从而 \(A=\kappa(h)\).
由于 \(\mathfrak n\mathcal O_{G,\xi}\) 包含在 \(\mathfrak m\) 中,
从而幂零, Nakayama 引理随即表明
\[
  \mathcal O_{H,h}\longrightarrow\mathcal O_{G,\xi}
\]
是满射.

设 \(V\) 是 \(H\) 的包含 \(h\) 的一个仿射开子集,
\(U\ne\varnothing\) 是 \(G\) 的包含在 \(u^{-1}(V)\)
中的一个仿射开子集, 并设
\[
  \varphi:\mathcal O_H(V)\longrightarrow\mathcal O_G(U)
\]
是 \(u\) 诱导的 \(k\)-代数同态.
设 \(\mathfrak p\) 是 \(\mathcal O_G(U)\) 中对应于
\(\xi\) 的素理想,\begingroup
\renewcommand{\thefootnote}{\fnsymbol{footnote}}%
\footnote[1]{译者注: 法文原文在此印为 \(\eta\);
紧接上文所引入的泛点是 \(\xi\).}%
\endgroup
并置 \(\mathfrak q=\varphi^{-1}(\mathfrak p)\).
由于 \(G\) 在 \(k\) 上有限型,
\(k\)-代数 \(\mathcal O_G(U)\) 由有限多个元素
\(a_1,\ldots,a_n\) 生成.
由以上论证, 存在元素 \(b_1,\ldots,b_n,s\in\mathcal O_H(V)\),
满足 \(s\notin\mathfrak q\), 使得在 \(\mathcal O_G(U)_{\mathfrak p}\) 中有
\[
  a_i/1=\varphi(b_i)/\varphi(s).
\]
因此存在元素
\[
  t_1,\ldots,t_n\in\mathcal O_G(U)-\mathfrak p
\]
使得
\[
  t_i\bigl(a_i\varphi(s)-\varphi(b_i)\bigr)=0.
\]
置
\[
  t=t_1\cdots t_n\varphi(s)\in\mathcal O_G(U)-\mathfrak p.
\]
则等式 \(a_i/1=\varphi(b_i)/\varphi(s)\)
已经在 \(\mathcal O_G(U)_t\) 中成立.
由于
\[
  t\in\mathcal O_G(U)=k[a_1,\ldots,a_n],
\]
存在一个元素 \(b\in\mathcal O_H(V)\),
使得对某个 \(r\in\mathbb N\), 有
\[
  t/1=\varphi(b)/\varphi(s)^r.
\]
所以 \(\varphi\) 诱导满射
\[
  \mathcal O_H(V)_{sb}\longrightarrow\mathcal O_G(U)_t,
\]
从而 \(u\) 在 \(\xi\) 处是局部浸入.

由那些 \(u\) 在其处为局部浸入的点所成的 \(G\) 的开子集 \(W\),
因而是非空的.
由于 \(G\) 是 Jacobson 概形, \(W\) 包含一个闭点 \(y\).
要证明 \(W=G\), 只需证明 \(G\) 的每个闭点都属于 \(W\).
然而, 每个闭点 \(x\) 都是 \(y\) 在平移
\(r_x\circ r_y^{-1}\) 下的像, 因而属于 \(W\).
因此 \(W=G\), 这证明 \(u\) 是局部浸入.

由于 \(G\) 不可约, 可知 \(u\) 是浸入.
事实上, 对每个 \(x\in G\),
设 \(U_x\) 和 \(V_x\) 分别为 \(G\) 与 \(H\) 的开子集,
使得 \(x\in U_x\), 且 \(u\) 诱导 \(U_x\) 到 \(V_x\) 的闭浸入.
由于 \(U_x\) 在 \(G\) 中稠密,
子集 \(u(U_x)\) 在 \(u(G)\cap V_x\) 中稠密.
但 \(u(U_x)\) 在 \(V_x\) 中是闭的, 所以
\footnote{Canon 来源说明 (SGA-SRC-000858): 重版印为 \(u(U_x)=V_x\).
前一句的稠密性陈述及下一条逆像恒等式确定了此处所写的交集.}
\[
  u(U_x)=u(G)\cap V_x.
\]
此外, 由于 \(u\) 是单射, \(U_x=u^{-1}(V_x)\).
因此, \(u\) 诱导 \(G\) 到 \(H\) 的由这些 \(V_x\) 覆盖的开子概形的闭浸入.
所以 \(u:G\to H\) 是浸入.
以上已知 \(u(G)\) 在 \(H\) 中是闭的, 因而 \(u\) 是闭浸入.

% Translation coverage: entire assigned component, retaining all source statements and proofs.
% Statistics: source lines=262; target lines=270; inline math=201; display math=23; unresolved translation TODOs=0; independent compilation not performed here.
